% !TeX root = ../main.tex
% ---------------------------------------------------------------------------
% Settlement, Building Damage and Protective Measures
% Scope extension: extends the Mair settlement material in Chapter 8 ( Chapter~\ref{ch:numerical-modelling})
% into third-party impact assessment.
% Key sources: Burland & Wroth (1974); Boscardin & Cording (1989); Mair et al.
% (1996) staged assessment; Potts & Addenbrooke (1997); CIRIA C760.
% ---------------------------------------------------------------------------
\chapter{Settlement, Building Damage and Protective Measures}
\label{ch:settlement-damage}


\section{The ground movement field}
\label{sec:the-ground-movement-field}

\subsection{From volume loss to damage}
\label{sec:volume-loss-to-damage}
\scopenote{The chain from construction control, through ground movement, to
building strain and damage category. Where in that chain design decisions
actually bite.}

\subsection{The three-dimensional settlement field}
\label{sec:3d-settlement}
\scopenote{Transverse Gaussian trough, longitudinal cumulative probability
profile, and the settlement history experienced by a point as the face passes.}

\subsection{Subsurface and horizontal movements}
\label{sec:subsurface-movements}
\scopenote{Variation of trough width with depth, horizontal displacement and
strain fields, and their importance for piled foundations and buried services.}

\section{Building damage assessment}
\label{sec:building-damage-assessment}

\subsection{Staged damage assessment}
\label{sec:staged-assessment}
\scopenote{Preliminary screening on maximum settlement and slope, second-stage
greenfield assessment, and detailed third-stage assessment for buildings that
do not pass.}

\subsection{Limiting tensile strain and damage categories}
\label{sec:limiting-strain}
\scopenote{Burland's deep and shallow beam idealisation, combining bending and
horizontal strain, Boscardin \& Cording's interaction chart, and the
negligible-to-severe damage category scale.}

\subsection{Building stiffness and soil-structure interaction}
\label{sec:building-stiffness}
\scopenote{Relative stiffness modification factors after Potts \&
Addenbrooke; why greenfield assessment is conservative and by how much;
treatment of facade openings and building geometry.}

\section{Foundations and buried assets}
\label{sec:foundations-and-buried-assets}

\subsection{Foundation type and response}
\label{sec:foundation-response}
\scopenote{Strip, raft and piled foundations; pile response to tunnelling
below or beside the toe; downdrag and load redistribution in pile groups.}

\subsection{Buried services and utilities}
\label{sec:utilities}
\scopenote{Jointed and continuous pipelines, allowable joint rotation and
strain, and the condition of legacy services as the governing uncertainty.}

\subsection{Track, road and railway response}
\label{sec:track-response}
\scopenote{Twist, cant and alignment tolerances for operating railways; the
tight movement limits that often govern urban tunnel design.}

\section{Existing tunnels and excavation interaction}
\label{sec:existing-tunnels-and-excavation-interaction}

\subsection{Existing tunnels and adjacent structures}
\label{sec:existing-tunnels}
\scopenote{Tunnelling beneath or beside an existing tunnel, ovalisation and
joint opening of the existing lining, and typical proximity criteria.}

\subsection{Twin tunnel and multiple excavation interaction}
\label{sec:twin-tunnel}
\scopenote{Superposition and its limits, modification factors for the second
drive, and the effect of pillar width and construction sequence.}

\section{Subsidence over voids and legacy workings}
\label{sec:subsidence-over-voids-and-legacy-workings}

\subsection{Subsidence over voids and legacy workings}
\label{sec:void-subsidence}
\scopenote{A mechanism distinct from tunnelling volume loss: chimney caving
and sinkhole development above shallow voids, abandoned mine workings and
karst cavities. Crown pillar stability, bulking during collapse, and the
migration height at which caving arrests. After Brady \& Brown Ch.~17,
adapted to the civil case of building over legacy workings.}

\subsection{Investigating and treating legacy voids}
\label{sec:legacy-voids}
\scopenote{Records research, probe drilling and geophysics to locate
workings; grout filling, capping and structural bridging as treatments; and
the residual risk that remains after treatment.}

\section{Protection and monitoring}
\label{sec:protection-and-monitoring}

\subsection{Protective measures}
\label{sec:protective-measures}
\scopenote{Compensation grouting, underpinning, structural strengthening,
in-ground barriers and jet grouted walls; selection by cost, disruption and
reliability.}

\subsection{Monitoring, trigger levels and response plans}
\label{sec:damage-monitoring}
\scopenote{Instrumentation of third-party assets, green-amber-red trigger
levels, contractual and stakeholder framing, and the response plan that must
exist before the trigger is reached.}
